\chapter{Related Works}\label{ch:related}

Dynamic bottleneck prediction connects manufacturing-system analysis with temporal forecasting and relational representation learning. To clarify the position of PHGST-DM, this chapter separates the definition of the prediction task from the choice of operational indicators, the evidence for delayed manufacturing interactions, and the modelling architectures that may exploit those interactions.

\section{Dynamic Bottleneck Prediction Concepts and Task Scope}\label{dynamic-bottleneck-prediction-concepts-and-task-scope}

A throughput bottleneck constrains system output and must be interpreted in the context of production relationships, rather than identified solely by the highest utilization or a local waiting queue. In their systematic review, Skoogh et al.~(2023) examined throughput-bottleneck detection methods and the different ways they are operationalized. The identity of a bottleneck may change with resource availability, material supply, and operating conditions. Lai et al.~(2021), for example, combined production data and factory-physics knowledge to detect shifting bottlenecks in complex, non-serial manufacturing systems.

Detection and prediction address different questions. Detection describes a constraint that is observable in the current or recorded system state, whereas prediction estimates a future constraint from information already available. Subramaniyan et al.~(2021) distinguish identification, diagnosis, prediction, and prescription within AI-assisted bottleneck analysis. Continually updating a detection result is therefore not equivalent to issuing an advance warning, and a correct forecast does not by itself demonstrate an effective production intervention.

Prediction methods have followed several complementary directions. Li et al.~(2011) investigated throughput bottleneck forecasting through time-series analysis, and Subramaniyan et al.~(2018) used machine active periods and an autoregressive integrated moving average (ARIMA) approach. Huang et al.~(2019) examined how future bottleneck forecasts could inform proactive task dispatching. More recent neural methods incorporate changing conditions and production structure. Liu et al.~(2023) constructed station supply--demand indicators and combined dual-attention LSTM forecasting with concept-drift detection and transfer learning. Ma et al.~(2023) modelled spatial and temporal information for high-variety shops with complex routings. Lai et al.~(2023) forecast blockage and starvation using station relations before inferring future bottlenecks.

These methods differ not only in architecture but also in the information they predict. A continuous bottleneck score, a future blocked-time value, and a discrete event report are distinct targets. An event-oriented report must identify a resource and, when an event is present, specify its onset and duration; the remaining duration is especially relevant for a constraint already active at the forecast origin. PHGST-DM uses this event-level formulation and additionally estimates observable process conditions and order completion progress.

\section{Bottleneck Indicators and Input Representation}\label{bottleneck-indicators-and-input-representation}

Predictive modelling requires operational indicators with explicit resource, time-window, and process meanings. Queue length and waiting time characterize accumulation; activity ratios describe resource utilization; blockage and starvation measurements reflect different constraints on available work. A systematic review by Skoogh et al.~(2023) groups detection evidence into queue states, process states, and their combinations. No single group is sufficient to characterize every type of manufacturing constraint.

Temporal persistence is another essential aspect of indicator design. Subramaniyan et al.~(2016) constructed a real-time average active-period measure from machine states and timestamps. Unlike an instantaneous activity value, a sustained active period retains information about the duration of an operating condition. This supports the broader principle that isolated changes and persistent constraints should not automatically receive the same interpretation.

Operational evidence may be combined through process rules or summarized by data-driven similarity. Lai et al.~(2021) incorporated blockage and starvation information together with manufacturing structure, while Liu et al.~(2023) exploited supply--demand relations when constructing predictive inputs. In a complementary direction, Subramaniyan et al.~(2020) clustered bottleneck-indicator time series using dynamic time warping and hierarchical clustering to identify similar production dynamics. K-means can likewise summarize recurrent multidimensional observations, although cluster meanings depend on feature representation and the chosen distance measure (Jain, 2010). A cluster of similar operating trajectories is \normalbold{not}, by itself, a bottleneck label or a validated causal explanation.

Manufacturing inputs also mix numerical measurements and categorical information. Separate feature projections can preserve meaningful differences before integration, as exemplified by variable-specific processing in the Temporal Fusion Transformer (Lim et al., 2021). For a multi-entity factory, the model must preserve both semantic feature groups and the identities of workstations, workers, handling resources, and buffers. These considerations motivate PHGST-DM's grouped resource-state embedding and the separation between rule-derived supervision and training-only temporal pattern prototypes.

\section{Delayed Cross-resource Effects in Production Systems}\label{delayed-cross-resource-effects-in-production-systems}

Finite buffers provide a manufacturing-specific reason why the consequence of a disturbance may be delayed. An upstream failure need not immediately starve a downstream resource while intermediate WIP remains available. Conversely, an upstream resource may continue processing after downstream interruption until the intermediate storage becomes full. The observable effect depends on the disruption location, intervening storage, and production relationships rather than solely on the state of the eventual bottleneck.

Gu et al.~(2013) examined hidden maintenance opportunities in discrete and complex production lines, showing how buffers can create periods in which interruptions do not immediately reduce output. Gu et al.~(2015) further studied the prediction of failure-induced starvation and blockage intervals at bottleneck machines when disruptions occur at other upstream or downstream resources, with simulation and manufacturing-case validation. These studies establish the operational relevance of the time separation between an initiating change and its later effect.

Zhang et al.~(2022) formulated \emph{Bottleneck Time Delay} in a smartphone digital-twin workshop as the interval between an operational disruption and the resulting production stoppage at the bottleneck workstation. The analysis relates this interval to disturbance location and intermediate buffer occupancy, while distinguishing the bottleneck impact from the later effect at the end of the line. This literature supports delayed transmission of disruption consequences, but delayed transmission should not be conflated with a change in bottleneck identity: an existing bottleneck can be affected by a distant disruption without migrating.

The resulting predictive challenge is to use evidence from related resources \emph{before} the target resource becomes visibly constrained. Earlier queue, material, transportation, and buffer observations may contain signals relevant to a later resource-level event. Such evidence motivates historical cross-resource modelling, but a forecast can exploit lagged associations without explicitly estimating physical propagation times. This separation is central to the interpretation of the delay-aware component developed in Chapter 3.

\section{Graph-based Spatiotemporal Modelling and Event Prediction}\label{graph-based-spatiotemporal-modelling-and-event-prediction}

Manufacturing resources are connected through process precedence, transport services, material transfer, and shared tasks. Graph neural networks offer methods for combining node states with relational information (Wu et al., 2021). Graph convolutional networks aggregate information according to graph connectivity (Kipf and Welling, 2017), while graph attention networks learn weights for neighbouring-node contributions (Veličković et al., 2018). Dynamic forecasting requires these relationships to be modelled together with temporal variation.

Traffic forecasting has supplied several foundational spatiotemporal architectures. DCRNN integrates diffusion-based spatial modelling and recurrence (Li et al., 2018); STGCN combines graph and temporal convolution (Yu et al., 2018); Graph WaveNet uses learned spatial dependencies (Wu et al., 2019); and STSGCN models synchronous spatial--temporal interactions (Song et al., 2020). These architectures demonstrate the utility of joint graph and history modelling. Nevertheless, road-network connectivity cannot be treated as equivalent to the process, handling, buffer, and labour relationships present in manufacturing.

Manufacturing graph predictors also vary in what their nodes and outputs represent. \normalbold{BSTAN} models production stations and their relations, combining graph attention and recurrent learning to predict blockage and starvation before identifying future bottlenecks (Lai et al., 2023). \normalbold{BTGCN} represents ERP-managed manufacturing tasks with links through shared operators; its graph--recurrent model forecasts task-execution frequency for bottleneck identification (Yin et al., 2025). \normalbold{Dual-BDSTN} combines temporal representations with dynamic, material-flow-related graph dependencies and cross-attention to predict blockage and starvation times (Cen et al., 2025). Their contributions show the value of manufacturing relationships while also making clear that workstation, task, and material-flow graphs are not interchangeable.

Attention-based forecasting extends the range of relationships that can be combined. The Transformer architecture uses self-attention for learning dependencies across representations (Vaswani et al., 2017), and its adaptations to forecasting have been reviewed by Wen et al.~(2023). GMAN combines spatial and temporal attention with historical-to-future transformation (Zheng et al., 2020). \normalbold{PDFormer} uses geographic and semantic attention masks and short-term pattern information in propagation-delay-aware traffic forecasting (Jiang et al., 2023). These concepts are relevant to PHGST-DM, but the factory graph must represent manufacturing-specific relations and its predictive attention weights must not be interpreted as confirmed causal paths.

A separate modelling question concerns the output target. Continuous forecasting predicts future measurements, whereas event prediction estimates occurrences and temporal attributes. Tax et al.~(2017) used LSTM-based process monitoring to predict subsequent activities, their timing, and remaining completion time. \normalbold{STGNPP} combined traffic-state representations, road relations, congestion history, a continuous-time GRU, and a point-process intensity model to predict congestion-event times and durations (Jin et al., 2023). These studies illustrate event-oriented output design, rather than implying that a manufacturing bottleneck has the same operational definition as traffic congestion. PHGST-DM draws on spatiotemporal representation and event-target concepts while constructing manufacturing-specific labels and multi-task outputs.

\section{Research Gaps}\label{research-gaps}

The reviewed research identifies three gaps that motivate the project. \normalbold{First, the observation scope remains incomplete for heterogeneous production systems.} Representative predictors organize workstations, machines, or tasks, but workers, handling devices, buffers, and material-supply evidence are not consistently represented together with their process--logistics relationships. Important early evidence can therefore fall outside the model's chosen graph.

\normalbold{Second, process-consistent supervision for sustained bottleneck events is difficult.} Scalar indicators and simple threshold rules may confuse brief fluctuations with persistent constraints and provide limited information about accompanying operating conditions. Event and cause targets should instead be traceable to observable production states and include explicit temporal-persistence criteria. Pattern clustering can summarize recurring behaviours, but should remain separate from event ground truth.

\normalbold{Third, delayed cross-resource evidence is not yet fully integrated into heterogeneous event prediction.} Manufacturing studies demonstrate that storage and resource relationships can delay the appearance of disruptions, while general spatiotemporal architectures offer ways to combine historical information. A manufacturing predictor still needs to organize earlier observations across resource types to estimate where and when future bottleneck events will occur. These gaps motivate, respectively, PHGST-DM's heterogeneous process--logistics graph, process-informed supervision and temporal pattern clustering, and recurrent-pattern-guided delay-aware modelling.
